\section{Introdução}

Este trabalho apresenta a implementação de um sistema de realidade aumentada (AR) utilizando marcadores fiduciais ArUco. O projeto foi desenvolvido como parte do exercício programa 1 (EP1) da disciplina MAC5915, envolvendo três componentes principais:

\begin{itemize}
    \item \textbf{Calibração de Câmera}: Determinação dos parâmetros intrínsecos da câmera através de um padrão de xadrez
    \item \textbf{Realidade Aumentada com Cubo}: Implementação básica renderizando um cubo 3D sobre marcadores ArUco
    \item \textbf{Realidade Aumentada com Bola de Tênis}: Extensão da implementação anterior com renderização de uma esfera texturizada
\end{itemize}

\subsection{Objetivos}

Os objetivos principais deste projeto são:

\begin{enumerate}
    \item Implementar um sistema de calibração de câmera utilizando o padrão de xadrez do OpenCV
    \item Desenvolver um sistema de tracking de marcadores ArUco em vídeo
    \item Integrar OpenCV com OpenGL para renderização de objetos 3D
    \item Implementar conversão adequada entre sistemas de coordenadas (OpenCV para OpenGL)
    \item Renderizar objetos virtuais texturizados sobre marcadores reais
\end{enumerate}

\subsection{Estrutura do Relatório}

Este relatório está organizado da seguinte forma: a Seção~\ref{sec:metodologia} descreve a metodologia utilizada em cada componente do projeto; a Seção~\ref{sec:resultados} apresenta os resultados obtidos; a Seção~\ref{sec:discussao} discute os resultados e limitações; e a Seção~\ref{sec:conclusao} apresenta as conclusões e trabalhos futuros.
